% [VERIFY] experiments trace to: experiments.md E01-E06; Tables 2,3,4,5,6,7,8 + derived; Figs 1,4,6,7; C01-C08
\section{Experiments}
\label{sec:exp}

We evaluate residual learning on ImageNet 2012 classification (1.28M training images, 50k
validation, 100k test, 1000 classes), on CIFAR-10 ($50$k training / $10$k test $32\times32$
images, 10 classes), and on PASCAL VOC and COCO object detection. Unless stated otherwise,
ImageNet results use 10-crop testing and all plain and residual variants share the single
training recipe of \cref{sec:method:impl} so that every comparison is controlled.

\subsection{Plain versus residual networks on ImageNet}
\label{sec:exp:imagenet}
We first compare 18- and 34-layer plain networks against their residual counterparts under
identical training, where the residual nets add only parameter-free \optA{} identity
shortcuts and therefore have no extra parameters (verifying the diagnosis and the
elimination of the \degprob). \Cref{tab:plainvsres} reports the controlled comparison. The
deeper \plainnet is \emph{worse}: plain-34 reaches 28.54\% top-1 error against 27.94\% for
plain-18. The training curves in the underlying analysis show that plain-34 sits above
plain-18 throughout optimization, not merely at convergence, so the gap is a degradation of
the optimization rather than overfitting; we further note that the deep \plainnet still
reaches competitive accuracy, indicating the solver makes progress only slowly. Adding
identity shortcuts reverses the ordering: \resnet-34 attains 25.03\% top-1, which is 3.51
points better than plain-34 and 2.85 points better than \resnet-18 (27.88\%). At 18 layers
the plain and residual nets reach comparable final accuracy (27.94 vs.\ 27.88), but the
residual variant converges faster from the start, indicating that residual learning eases
optimization even where the plain net can ultimately converge. Together these results
support both the existence of the \degprob and its removal by residual learning.

% Source: evidence/tables/table2_imagenet_plain_vs_residual.md (Table 2). Supports C01, C02.
% EXACT numbers, 10-crop top-1 val error %.
\begin{table}[t]
\centering
\caption{Top-1 error (\%, 10-crop testing) on ImageNet validation. The residual nets here
have no extra parameter compared to their plain counterparts (\optA{} identity
shortcuts). The deeper \plainnet degrades (34 worse than 18); the deeper residual net
improves (34 better than 18). Source: Table~2.}
\label{tab:plainvsres}
\small
\begin{tabular}{l c c}
\toprule
 & plain & \resnet \\
\midrule
18 layers & 27.94 & 27.88 \\
34 layers & 28.54 & \best{25.03} \\
\bottomrule
\end{tabular}
\end{table}


\begin{figure}[t]
\centering
\includegraphics[width=0.6\linewidth]{fig_imagenet_plain_vs_resnet.pdf}
\caption{ImageNet top-1 validation error (\%, 10-crop) for plain versus residual networks
at 18 and 34 layers (\optA{} identity shortcuts, no extra parameters). The deeper
\plainnet degrades ($27.94\rightarrow28.54$) while the deeper residual net improves
(\textbf{$27.88\rightarrow25.03$}), reversing the ordering. Bars are the exact Table~2
values. $N=1$ trained model per bar.}
\label{fig:plainvsres}
\end{figure}

\subsection{Identity versus projection shortcuts}
\label{sec:exp:shortcut}
We next ask whether the parameter-bearing projection shortcuts of \cref{eq:projblock} are
necessary, ablating the three shortcut options on \resnet-34 under the same recipe.
\Cref{tab:shortcut} isolates the relevant rows. All three options improve over plain-34 by
a wide margin (3.51 to 4.35 top-1 points), confirming that the gain comes from residual
learning rather than from a particular shortcut form. The differences among the options are
small: \optA{} (identity, parameter-free) reaches 25.03\%, \optB{} (projection only where
dimensions change) reaches 24.52\%, and \optC{} (projection on every shortcut) reaches
24.19\%. The slight B-over-A margin is attributable to the zero-padded extra channels in A
carrying no residual learning, and the slight C-over-B margin to the parameters added by the
13 extra projection shortcuts. Because \optC{} is only marginally better yet costs extra
parameters and memory, it is not essential for addressing the \degprob; we therefore default
to identity shortcuts and reserve projections for dimension-matching, which is especially
important for the bottleneck designs where a projection on the high-dimensional ends would
roughly double a block's complexity.
% EXP-MORE

\subsection{Going deeper with bottleneck networks}
\label{sec:exp:depth}
Using the bottleneck block of \cref{sec:method:arch} with \optB{} shortcuts, we scale depth
to 50, 101, and 152 layers and evaluate whether accuracy continues to improve with depth
(verifying the depth-scaling and bottleneck claims). \Cref{tab:imagenetfull} reports the
full 10-crop validation table. Top-1 error decreases monotonically with depth---22.85\%
(\resnet-50), 21.75\% (\resnet-101), 21.43\% (\resnet-152)---and every residual variant
improves substantially over the 34-layer \plainnet (28.54\%) and over the cross-family
baselines VGG-16 (28.07 / 9.33) and PReLU-net (24.27 / 7.38). The gains come from depth
rather than added parameters relative to the plain design, and from \cref{tab:arch} the
152-layer model has $11.3\times10^{9}$ FLOPs, still below VGG-19's $19.6\times10^{9}$
despite being eight times deeper. \Cref{fig:resnetdepth} plots the monotone trend for both
metrics. Under the stronger single-model multi-scale fully-convolutional protocol
(\cref{tab:singlemodel}), \resnet-152 reaches 19.38 / 4.49 top-1/top-5, surpassing every
prior single model including PReLU-net (21.59 / 5.71) and BN-inception (21.99 / 5.81).
Finally, a six-model ensemble attains 3.57\% top-5 error on the ImageNet \emph{test} set
(\cref{tab:ensemble}), improving on the best prior ensemble (BN-inception, 4.82\%) and
winning the first place in ILSVRC 2015 classification.

% Source: evidence/tables/table3_imagenet_validation_full.md (Table 3). Supports C01,C02,C03,C04,C05.
% EXACT numbers, 10-crop val error %. ResNet-50/101/152 use option B (projections only on
% increasing dimensions); VGG-16 is the artifact's own test; GoogLeNet top-1 not reported.
\begin{table}[t]
\centering
\caption{Error rates (\%, 10-crop testing) on ImageNet validation. \resnet-50/101/152 use
\optB{} (projections only for increasing dimensions). Top-1 error decreases monotonically
with depth from 50 to 152, and every residual variant improves over the 34-layer
\plainnet. Source: Table~3.}
\label{tab:imagenetfull}
\small
\begin{tabular}{l c c}
\toprule
model & top-1 err. & top-5 err. \\
\midrule
VGG-16~\citep{simonyan2015vgg}      & 28.07 & 9.33 \\
GoogLeNet~\citep{szegedy2015googlenet} & ---   & 9.15 \\
PReLU-net~\citep{he2015prelu}       & 24.27 & 7.38 \\
\midrule
plain-34          & 28.54 & 10.02 \\
\resnet-34 A      & 25.03 & 7.76 \\
\resnet-34 B      & 24.52 & 7.46 \\
\resnet-34 C      & 24.19 & 7.40 \\
\resnet-50        & 22.85 & 6.71 \\
\resnet-101       & 21.75 & 6.05 \\
\resnet-152       & \best{21.43} & \best{5.71} \\
\bottomrule
\end{tabular}
\end{table}

% Source: evidence/tables/table4_imagenet_singlemodel.md (Table 4). Supports C03.
% EXACT numbers. Single-model, fully-convolutional multi-scale testing. dagger = test set.
\begin{table}[t]
\centering
\caption{Single-model error rates (\%) on ImageNet validation (multi-scale
fully-convolutional testing); $\dagger$ denotes a result on the test set. The
single-model \resnet-152 surpasses every prior single model. Source: Table~4.}
\label{tab:singlemodel}
\small
\begin{tabular}{l c c}
\toprule
method & top-1 err. & top-5 err. \\
\midrule
VGG~\citep{simonyan2015vgg} (ILSVRC'14)        & ---   & 8.43$^\dagger$ \\
GoogLeNet~\citep{szegedy2015googlenet} (ILSVRC'14) & --- & 7.89 \\
VGG~\citep{simonyan2015vgg} (v5)               & 24.4  & 7.1 \\
PReLU-net~\citep{he2015prelu}                  & 21.59 & 5.71 \\
BN-inception~\citep{ioffe2015batchnorm}        & 21.99 & 5.81 \\
\resnet-34 B    & 21.84 & 5.71 \\
\resnet-34 C    & 21.53 & 5.60 \\
\resnet-50      & 20.74 & 5.25 \\
\resnet-101     & 19.87 & 4.60 \\
\resnet-152     & \best{19.38} & \best{4.49} \\
\bottomrule
\end{tabular}
\end{table}

% Source: evidence/tables/table5_imagenet_ensembles.md (Table 5). Supports C03.
% EXACT numbers, top-5 error % on the ImageNet TEST set (reported by the test server).
\begin{table}[t]
\centering
\caption{Ensemble error rates (\%, top-5) on the ImageNet \emph{test} set, reported by the
test server. The six-model \resnet ensemble attains 3.57\%, winning first place in
ILSVRC 2015 classification. Source: Table~5.}
\label{tab:ensemble}
\small
\begin{tabular}{l c}
\toprule
method & top-5 err.\ (test) \\
\midrule
VGG~\citep{simonyan2015vgg} (ILSVRC'14)        & 7.32 \\
GoogLeNet~\citep{szegedy2015googlenet} (ILSVRC'14) & 6.66 \\
VGG~\citep{simonyan2015vgg} (v5)               & 6.8 \\
PReLU-net~\citep{he2015prelu}                  & 4.94 \\
BN-inception~\citep{ioffe2015batchnorm}        & 4.82 \\
\resnet (ILSVRC'15)                            & \best{3.57} \\
\bottomrule
\end{tabular}
\end{table}


\begin{figure}[t]
\centering
\includegraphics[width=0.62\linewidth]{fig_resnet_imagenet_depth.pdf}
\caption{ImageNet validation error (\%, 10-crop) versus residual depth; \resnet-34 uses
\optA{} and the deeper nets use \optB. Both top-1 and top-5 error fall monotonically as
depth grows from 34 to 152 (\textbf{top-1 $25.03\rightarrow21.43$}). Markers are the exact
Table~3 values; no curve is interpolated. $N=1$ trained model per depth.}
\label{fig:resnetdepth}
\end{figure}

\subsection{Analysis on CIFAR-10}
\label{sec:exp:cifar}
To test whether residual learning generalizes beyond ImageNet and to probe extreme depth,
we train the minimal $6n{+}2$-layer CIFAR family at depths 20, 32, 44, 56, 110, and 1202,
with plain nets of matched depth as controls (verifying the extreme-depth generalization and
warmup claims). \Cref{tab:cifar} and \cref{fig:cifar} report test error. Residual error
falls steadily with depth, from 8.75\% at 20 layers to 6.43\% at 110 layers (best; mean
$6.61\pm0.16$ over five runs), while the matched plain nets degrade---plain-110 exceeds 60\%
error and is omitted from the comparison. The 110-layer residual net required a learning-rate
warmup to begin converging cleanly: an initial rate of 0.1 was slightly too large at this
depth, so we warmed up at 0.01 for roughly 400 iterations until training error fell below
about 80\%, then restored 0.1; this is a stability aid for very deep CIFAR nets rather than a
fundamental requirement of residual learning. The 1202-layer net trains with no optimization
difficulty (final training error below 0.1\%), but its test error rises to 7.93\%, higher than
the 110-layer net. Because optimization is not the bottleneck at 1202 layers, we attribute
this to overfitting: a 19.4M-parameter model on a 50k-image dataset, trained here without
strong regularization such as maxout or dropout.

% Source: evidence/tables/table6_cifar10.md (Table 6). Supports C06, C07.
% EXACT numbers, CIFAR-10 test error %. ResNet-110 shows best (mean +/- std) over 5 runs.
\begin{table}[t]
\centering
\caption{Classification error (\%) on the CIFAR-10 test set; all methods use data
augmentation. For \resnet-110 we report best (mean $\pm$ std) over five runs. Error falls
as residual depth grows from 20 to 110; the 1202-layer net trains successfully but
overfits. Source: Table~6.}
\label{tab:cifar}
\small
\begin{tabular}{l c c c}
\toprule
method & \# layers & \# params & error (\%) \\
\midrule
Maxout~\citep{goodfellow2013maxout}     & ---  & ---   & 9.38 \\
NIN~\citep{lin2013nin}                  & ---  & ---   & 8.81 \\
DSN~\citep{lee2014dsn}                   & ---  & ---   & 8.22 \\
\midrule
FitNet~\citep{romero2015fitnets}        & 19   & 2.5M  & 8.39 \\
Highway~\citep{srivastava2015highway,srivastava2015trainingverydeep} & 19 & 2.3M & 7.54 (7.72 $\pm$ 0.16) \\
Highway~\citep{srivastava2015highway,srivastava2015trainingverydeep} & 32 & 1.25M & 8.80 \\
\midrule
\resnet & 20   & 0.27M & 8.75 \\
\resnet & 32   & 0.46M & 7.51 \\
\resnet & 44   & 0.66M & 7.17 \\
\resnet & 56   & 0.85M & 6.97 \\
\resnet & 110  & 1.7M  & \best{6.43 (6.61 $\pm$ 0.16)} \\
\resnet & 1202 & 19.4M & 7.93 \\
\bottomrule
\end{tabular}
\end{table}


\begin{figure}[t]
\centering
\includegraphics[width=0.62\linewidth]{fig_cifar10_depth.pdf}
\caption{CIFAR-10 test error (\%) versus residual depth (log scale; \optA{} identity
shortcuts). Error decreases monotonically through 110 layers
(\textbf{$8.75\rightarrow6.43$}) and then rises at 1202 layers (highlighted), where the
network still trains cleanly but overfits the small dataset. Markers are the exact Table~6
values; the 110-layer point is the best of five runs.}
\label{fig:cifar}
\end{figure}

\subsection{Object-detection transfer}
\label{sec:exp:detection}
Finally, we test whether deeper residual representations transfer to localization-sensitive
tasks by replacing the VGG-16 backbone of a baseline Faster R-CNN detector with \resnet-101,
holding the detection framework and hyperparameters fixed (verifying the transfer claim).
\Cref{tab:detection} reports the swap. On PASCAL VOC the mAP@.5 rises from 73.2 to 76.4 on
VOC07 and from 70.4 to 73.8 on VOC12. On COCO the mAP@.5 improves from 41.5 to 48.4, and the
stricter mAP@[.5,.95]---which rewards tighter localization---improves from 21.2 to 27.2, a 6.0
point absolute and 28\% relative gain. That the mAP@[.5,.95] gain is comparable to the
mAP@.5 gain indicates the deeper features help both recognition and localization, and because
only the backbone changed, the improvement is attributable to the learned representation.

% Source: evidence/tables/table7_pascal_voc_detection.md + table8_coco_detection.md (Tables 7,8). Supports C08.
% EXACT numbers, baseline Faster R-CNN, mAP %. VGG-16 vs ResNet-101 backbone swap.
\begin{table}[t]
\centering
\caption{Object-detection transfer with baseline Faster R-CNN: swapping a VGG-16 backbone
for \resnet-101 with otherwise identical hyperparameters. PASCAL VOC reports mAP@.5 (VOC07
test from 07+12 training; VOC12 test from 07++12 training); COCO reports mAP@.5 and
mAP@[.5,.95] on the validation set. The COCO mAP@[.5,.95] gain of 6.0 points is a 28\%
relative improvement. Source: Tables~7 and 8.}
\label{tab:detection}
\small
\begin{tabular}{l c c c c}
\toprule
 & \multicolumn{2}{c}{PASCAL VOC (mAP@.5)} & \multicolumn{2}{c}{COCO val} \\
\cmidrule(lr){2-3} \cmidrule(lr){4-5}
backbone & VOC07 test & VOC12 test & mAP@.5 & mAP@[.5,.95] \\
\midrule
VGG-16        & 73.2 & 70.4 & 41.5 & 21.2 \\
\resnet-101   & \best{76.4} & \best{73.8} & \best{48.4} & \best{27.2} \\
\bottomrule
\end{tabular}
\end{table}


% Source: evidence/tables/derived_from_table3_shortcut_options.md (derived from Table 3). Supports C04.
% EXACT numbers. Delta-vs-plain-34 carried from the derived subset.
\begin{table}[t]
\centering
\caption{\resnet-34 shortcut-option ablation on ImageNet validation (\%, 10-crop),
isolated from Table~3 with the plain-34 baseline. All three options beat plain-34 by a
wide margin; the differences among A, B, and C are small ($\le0.84$ top-1). Source:
derived subset of Table~3.}
\label{tab:shortcut}
\small
\begin{tabular}{l l c c c}
\toprule
model & shortcut option & top-1 err. & top-5 err. & $\Delta$ top-1 vs.\ plain-34 \\
\midrule
plain-34   & none                          & 28.54 & 10.02 & --- \\
\resnet-34 A & identity $+$ zero-pad        & 25.03 & 7.76 & $-3.51$ \\
\resnet-34 B & projection on dim.\ change   & 24.52 & 7.46 & $-4.02$ \\
\resnet-34 C & projection on every shortcut & 24.19 & 7.40 & $-4.35$ \\
\bottomrule
\end{tabular}
\end{table}


